\section{实验}
\label{sec:experiments}
\subsection{平台与任务}
本文在 Gazebo 物理仿真中评价固定系统 \texttt{paper1-eval-finite-scan-v1}。
P450 提供空中 RGB-D 与 MID360 观测，BUNKER 搭载 AUBO i5 机械臂、AG95 夹爪和腕部 D435。
公共软件栈使用 PX4 SITL/Prometheus 控制无人机、ROS 导航控制底盘，并使用基于 MoveIt 的机械臂规划。
每个任务从正常空中目标感知开始，通过策略自身的观测、站位选择和机器人执行推进，
不人工指定成功底盘位姿。目标是位于静态水平地面的已知砖块。
抓取通过仿真接触与正常控制器执行，不使用物体瞬移或替换成功标签。

所有方法从相同 map 初始观测位姿 $(x,y,z,\psi)=(-1.4,0,1.2,0)$ 开始，
使用相同 RGB-D 接受门限、对齐的 $0.10\,\mathrm{m}$ 场及适用的第~\ref{sec:method} 节参数。
实际采集版本分别为 AGENT \texttt{3d13a95}、SIM \texttt{a0ae8e3}、原始 RM4D \texttt{e9d4312}，
并使用独立校准的地面任务资产。这些标识对应真实执行版本，而不是后续开发分支的最新版。
采集期间未改变机器人决策规则、场景种子、传感器门限或物理成功判据。

\subsection{独立场景与配对协议}
\label{sec:scenes}
主研究包含 96 个独立抽样场景。Easy、Moderate、Hard 的抽样概率各为 $1/3$，
前瞻冻结后的实际数量为 38、28、30。
三档分别包含零、一、两道墙；墙高 $1\,\mathrm{m}$，中等场景墙长 $1\,\mathrm{m}$，
困难场景墙长 $0.6\,\mathrm{m}$，厚度均为 $0.15\,\mathrm{m}$。
目标位置/偏航、初始底盘位姿、墙位置/偏航按冻结规则独立扰动。
目标中心在 $(2,0)$ 附近，底盘在 $(3,-2.5)\,\mathrm{m}$ 附近。
目标 XY 扰动界为 $0.15\,\mathrm{m}$，偏航界为 $\pi/6$；
墙中心扰动界为 $0.15\,\mathrm{m}$，偏航界为 $\pi/12$。

Hard 生成规则引入遮挡机会，并保留不一定与操作相关的未知空间，
但不保证每个实例都产生有用区域阻断或有利的方法差异。
种子与开发数据互不重叠，完整几何、方法顺序及辅助子集均在运行方法前固定。
可接受性检查只涉及初始几何和传感器设置，不依据候选数量、路径存在性、视点评分或取回结果。
每个场景按预定平衡/打乱顺序分别运行两个主方法，共 192 次计划主任务。
独立分析单位是配对场景，而非视点或候选。分析保留实际抽样难度混合，
不在观察结果后重新按难度加权。

\subsection{对照与消融}
\label{sec:baselines}
唯一主对比是 Generic NBV 与 Ours。Generic 使用式~\eqref{eq:genericgain}，
Ours 使用式~\eqref{eq:taskgain}。两者共用初始条件、候选生成、有限扫描机会与遮挡、
成本、观测预算、证据更新、精确支撑、操作门控、候选筛选及完整执行链。
后续测量可因动作不同而不同。Generic 同时就是去除任务加权的消融，不重复计为第二组实验。

两种背景对照在六个预定场景上运行，即每档最先生成的两个编号。
“仅 RM4D”使用初始空中 RGB-D 目标和校准后的任务域资产，选择 baseline 首位候选，
跳过 MID360 确认及交接前多候选筛选，但保留下游物理检查。
因此，它是背景对照，而不是对相关性加权的严格隔离或对原文献地图工作域的测试。
“固定视点”在当前实测悬停位置重复观测，不选择平移 NBV，保留三个窗口上限及共用筛选停止条件。

在最先生成的四个 Hard 编号上，“去除遮挡预测”仅从预测增益中移除遮挡，不移除碰撞阻断或执行检查；
“去除飞行成本”将策略成本权重置零。各辅助比较复用预定的 Ours 对应任务。
12 次背景对照与 8 次消融使总任务数达到 212，而非 212 次主实验。
这些小子集用于描述性解释，不支持有统计功效的次要优势、等价性或普遍必要性主张。

\subsection{终点、失败与观测成本}
\label{sec:metrics}
主二元终点为物理取回成功，要求感知、到位、抓取、提升及短时保持检查全部满足，
其中包含独立仿真物体高度检查。现有检查器要求 TCP 和物体高度均至少增加
$0.10\,\mathrm{m}$；任务同时要求持续一秒墙钟时间的最新抓取确认。
该接口保持与下文的仿真时间效率指标不同。
候选确认、规划预览成功、$\Dexec$、TCP 运动或任务进程的成功消息，单独均不等于取回成功。
记录阶段包括观测、确认、筛选、导航/停车、D435 精定位、$\Dexec$、抓取、提升/保持。
首次终止分类为每个失败任务分配一个原因；未到达的下游阶段不额外计为失败。
模型拒绝也不被解释为已真实发生物理碰撞。

共用预算为三个被接受的五仿真秒观测窗口，包含初始窗口。
追加观测可以发生在平移、仅偏航或基本相同的位置重观测之后。
本文区分窗口数与实测观测位姿之间真正的位置变化。
平移超过 $0.20\,\mathrm{m}$ 记为可判定位置变化；包裹后的偏航差超过
$0.20\,\mathrm{rad}$ 可单独识别。这些是报告分类，不是路径长度或新增控制门限。

机器人效率时间均使用仿真时间。主动感知时间从首次采集开始，至主动停止或终止失败结束，
包含观测、运动、等待、计算和共用预览，但不包含初始起飞、空中 RGB-D 和首次可达性查询。
地面执行时间从导航开始到地面任务终止事件结束。
全任务时间包含初始化及正常返航降落。未到达地面阶段时，其时间保持缺测，不记为零。
该版本距离轨迹记录不完整，因此不把距离节省作为推断终点或论文主张。

正常观测不足、有界规划拒绝与物理执行失败均为有效结果。
只有有记录证明的基础设施无效尝试，才能在固定储备内进行一次同场景、同方法替代。
全部 212 个计划 slot 均有有效的最终二元结果，96 组主配对完整。
两个确认无效的原始尝试获得替代，形成 214 次正式激活；另一次外部裸仿真启动独立计入资源。
原始记录均保留。这些激活不构成额外独立样本。

\subsection{预先规定的分析}
\label{sec:analysis}
令 $a,b,c,d$ 分别表示两者成功、仅 Ours 成功、仅 Generic 成功、两者失败的场景数。
报告边际成功率以及
\begin{equation}
\widehat\Delta=(b-c)/96.
\label{eq:effect}
\end{equation}
唯一主检验是以 $b+c$ 个不一致配对为条件的双侧精确 McNemar 检验，显著性水平为 $0.05$。
预定保守 95\% 差值区间由 $b/96$、$c/96$ 各自的 97.5\% Clopper--Pearson 区间
$[L_b,U_b]$、$[L_c,U_c]$ 组合为 $[L_b-U_c,U_b-L_c]$。
Bonferroni 联合覆盖率至少为 95\%，不假定两个不一致类别独立。
该区间针对配对风险差，而非 $p$ 值。
首次激活敏感性分析将基础设施无效原始任务计为未成功，不删除分母。

难度分层和辅助比较均为描述性分析，不是额外主检验。
条件效率在 53 个两种方法均成功的场景上汇总，按整个配对场景进行 10,000 次 bootstrap 重采样，
采用固定分析种子。条件性需要明确：这些均值不是无偏的全任务成本效应，
快速失败也不是高效取回。已有分析同时提供分母为全部 96 场景的成功—资源曲线。
本文写作不引入新实验、新显著性检验或结果驱动的样本扩展。
